\begin{frame}[t]{\ev{\tagP}Test 1 compares restore with a warm cached template}
\lead{Pre-registered and unrun. No model calls.}
{\footnotesize\renewcommand{\arraystretch}{1.0}
\begin{tabularx}{\textwidth}{@{}L{3.2cm}Y@{}}
\toprule
Arm & Starting point\\
\midrule
Warm template & dependencies baked in, source fetched at the pinned commit\\
Restore & captured reached state after setup and one warm test run\\
Cold reference & no cache, used as a reference only\\
\midrule
Checks first & record whether a RAM nonce survives (else report fan-out), and restored and direct outputs match\\
Design & $N\in\{3,6,12\}$ concurrent cells, 20 interleaved batches per arm and $N$\\
Outcome & time from the batch request until the last child's first test result\\
Supports & lower bound of the 95\% CI of the median gain above 10\,s at every $N$\\
Rejects & upper bound of the 98.3\% (Bonferroni) CI below 10\,s at any $N$\\
\bottomrule
\end{tabularx}\par}
\claim{A result in which the cache wins is a valid outcome.}
\sources{Pre-registration, tags prereg-v1 to prereg-v3. Restore failures are reported with the timing results.}
\note{[0:55] Section five, evaluation, presents two pre-registered tests. Neither has run. Test one compares restore with a warm cached template. First we record whether a value held in memory survives restore. Restored outputs must also match direct outputs. We then run three, six and twelve cells, twenty batches per arm and size. The outcome is the time until the last child reports its first test result. Restore is supported only if the lower confidence bound of the gain exceeds ten seconds at every size. A cache win is a valid outcome. The second test concerns failures.}
\end{frame}
